\section{Concluding Remarks}
\label{sec:concl}

\paragraph{Discussions and Future Work}
In the future, we aim to expand our framework's capabilities and address its current limitations. First, our implemented symbolic executor currently lacks support for recursive paging functions; we plan to lift this limitation to enable verification of more complex page table management implementations in Linux. Second, our framework currently requires the manual extraction of paging functions from the target system. This process is error-prone, particularly when parsing intricate header files and C macros within large codebases. To resolve this, we plan to develop automated tooling to identify and extract the necessary paging functions directly from complex systems. Finally, our machine model does not account for TLBs and caches. However, prior efforts~\cite{certikos,li:ieeesp2021,taososp21,li:sekvm:usenix21} demonstrate that system code managing TLBs, caches, and page table synchronization can be verified independently using a layered, modular approach that is independent of proofs for the paging function. Our exclusive focus on verifying core paging functions remains orthogonal to, and fully composable with, these existing efforts. Spoq2~\cite{spoq2} was recently introduced to automate verification of security properties in systems software by associating implementations to specifications first and then verifying that these specifications satisfy the intended properties. We will explore integrating PPV with Spoq2 to facilitate security reasoning of page table management code by utilizing an array-like representation.

\paragraph{Conclusions.} Verifying the correctness of low-level system code, particularly functions managing page tables, demands significant proof effort. Although prior work employing structural lifting helped reduce this cost, it still required substantial manual labor. To overcome this challenge, we introduced a novel, push-button approach that scales verification by encoding all proof tasks as quantifier-free formulas, thereby enabling their full automation through efficient SMT solvers. We successfully applied this methodology to prove multiple paging function implementations adapted from a Linux KVM hypervisor, demonstrating their correct operation with multi-level page tables and support for essential features, including dynamic allocation and huge pages. Furthermore, we demonstrated that our framework integrates seamlessly with existing full-system verification tools. Our results validate the feasibility and efficiency of using a push-button approach to practically verify paging functions from commodity systems.
